% Included by dominated_rho8.tex. Requires its notation and theorem environments.
\section{The calibrated implementation with the original common input}
\label{sec:dominated-common-input}

The implementation considered here retains the original common additive
read and all three non-AX branches.  The theorem proved in this section concerns
that implementation, including the source rebracketing ultimately used
for the proposed machine.  Preserving source addresses is useful for
retaining an established initializer, but is not by itself a proof of a
binary blank-tape startup.

\subsection{Exact command semantics and macro evaluation}

Write the retained tuple as $(t,w,d,h)$ before the common read.  The
stack notation $[x_1,\ldots,x_k;w]$ denotes right-nested pairing,
with $\unp(0)=(0,0)$.  The exact non-AX command equations are
\begin{align*}
 \MP(z)(t,w,d,h)&=
 \begin{cases}
 (d+z,\pr(0,w),d+z,0),&I(t,d+z)=h,\\
 (0,w,d+z,0),&I(t,d+z)\ne h;
 \end{cases}\\
 \DROP(z)(t,w,d,h)&=(h,\pr(t,w),d+z,h),\\
 \NIP(z)(t,w,d,h)&=
 \begin{cases}
 \mathsf{halt},&h=1,\\
 (t,w',d+z,h'),&(h',w')=\unp(w),\ h\ne1.
 \end{cases}
\end{align*}
Even DROP and NIP retain an additive field; all compiler uses of their
fields below are zero.
The AX arguments are written in actual wire order as
\[
 \AX(c,z,y,x),\qquad
 f(\AX(c,z,y,x))=[c+3,z,y,x].
\]
Its numerical parameters are
\[
 U=d+z,\qquad A=U+x,\qquad B=h+y,\qquad C=c.
\]
Define the gated top by
\[
 G(A,B,c)=\begin{cases}
 0,&c=A\text{ or }c=B,\\
 \operatorname{bundle}_{219}(A,B,c),&\text{otherwise}.
 \end{cases}
\]
Write $R_1,\ldots,R_7$ for the seven retained rows, including the merged
equality row, in the inherited stack order, and write
$Q_Uh=4\pr(U,h)+2$ for the universal-quantifier code.  The AX output is
\begin{align*}
 t'&=G(A,B,c),&d'&=0,&h'&=0,\\
 w'&=[R_1(A,B,c),\ldots,R_7(A,B,c),L(A,B,c,U),Q_Uh,t;w].
\end{align*}
Here the body of the generalized formula is the \emph{incoming} cache
$h$, before addition of $y$.  Since $U\leq U+x=A$, every raw AX command
produces an admissible dominated-row instance.

The reference source \texttt{a194\_b194\_f1.json} keeps the old common
read before dispatch.  Its replaced AX prefix has the closed macro word
\[
 \begin{aligned}
 &a,\ b,\ Q,\ c,\ a,\ \operatorname{read}_H,\
       \operatorname{read}_D,\ I,\\
 &a,\ I,\ a,\ c,\ I,\ b,\ c,\ I,\ I,\ I.
 \end{aligned}
\]
The macro $a$ loads the current $D$ while preserving it and pushing the
old top; $b$ does the same for $H$; and $c$ loads the selector.
After $a,b,Q$, the formula $Q_Uh$ has been constructed.  The next $c,a$
places $C,U$ in the partially constructed axiom.  The two field readers
then replace $H$ by $B$ and $D$ by $A$ without changing the loaded top
$U$.  The next $I$ forms $C\Imp U$.  All later $a$ loads read $A$ and
all later $b$ loads read $B$, so the remaining word constructs exactly
$L(A,B,C,U)$.  The inherited subsequent construction of the other rows
and the gated bundle is unchanged.  Every component here is an inherited
all-natural, normally returning macro, restoring zero scratch registers
at its boundary.

The theoremhood invariant is therefore preserved: top, cache, and
every stack head are dominated-calculus theorems.  MP applies the exact
numerical modus-ponens rule on success, and inserts truth on failure.
DROP and nonhalting NIP only move theorem codes; their additive counters
are not assumed to be theorem codes.  Generalizing
an arbitrary cached theorem with the arbitrary natural binder $U$ is
allowed.  A halt can occur only when a NIP finds theorem cache $1$.
This establishes command soundness for every raw candidate, including
commands with nonzero DROP or NIP additive fields.

\subsection{An explicit strong block compiler}

A block for $F$ must halt or return $(F,w,0,0)$ from every $(t,w,0,0)$.
The following blocks use only the exact command equations above.
Define
\[
 D_0=[\MP(1),\AX(1,0,0,0)]\;\NIP(0)^{10}.
\]
It returns $(0,w,0,0)$ from every $(t,w,0,0)$: MP fails because
$I(t,1)>0$, then AX has $U=A=C=1$ and zero gate.  Ten NIPs discard its
nine generated entries and saved zero top.  Those nine entries are
nonfalsity codes, so no halting test interrupts this calculation.
Thus $D_0$ is also the truth block.

For a dominated-row instance $L(a,b,c,u)$ with $u\leq a$, use
\[
 D_0\;[\AX(c,u,b,a-u)]\;\NIP(0)^8\;[\DROP(0)]\;\NIP(0)^3.
\]
The AX receives $d=h=0$, hence its parameters are exactly
$U=u$, $A=a$, $B=b$, $C=c$.  The suffix selects the eighth stack row,
then restores the underlying stack and zero cache.
For a retained row $R_j(a,b,c)$, $1\leq j\leq7$, use
\[
 D_0\;[\AX(c,0,b,a)]\;\NIP(0)^j\;[\DROP(0)]\;\NIP(0)^{11-j}.
\]
For a permitted gated bundle use
\[
 D_0\;[\AX(c,0,b,a)]\;\NIP(0)^{10}.
\]
In these last two cases the unused fourth argument is $U=0$, so it
satisfies the required bound for every $a$.

If $\mathcal B_F$ produces $F$, generalization with binder $i$ is compiled
as
\begin{multline*}
 \mathcal B_F\;[\DROP(0),\NIP(0),\AX(i,i,0,0)]\\
 {}\NIP(0)^9\;[\DROP(0),\NIP(0),\NIP(0)].
\end{multline*}
After the first two commands, the retained tuple is $(0,w,0,F)$.
The common AX field sets $U=i$, and $x=y=0$ gives $A=C=i$,
$B=F$.  Its gate is zero and its ninth stack entry is $Q_iF$.
The selection suffix restores $(Q_iF,w,0,0)$.

For modus ponens, given blocks for $F$ and $F\Imp G$, use
\begin{multline*}
 \mathcal B_F\;[\DROP(0)]\;\mathcal B_{F\Imp G}\;\\
 [\NIP(0),\DROP(0),\NIP(0),\MP(G),\AX(G,0,0,0)]\;\\
 \NIP(0)^{10}\;[\DROP(0),\NIP(0),\NIP(0)].
\end{multline*}
Unless a prior NIP halts, the two hypotheses reach top and cache in the
required order and MP succeeds.  The reset AX has $U=A=C=G$ and zero
gate; it preserves the derived conclusion among its saved stack entries.
The remaining commands return $(G,w,0,0)$ or halt if a cache already
contains $1$.  Early halting is permitted by the block contract.

Induction on a finite dominated-calculus derivation now proves this
strong block contract for every derived natural code.  No numerical
bound is tested at runtime: each row leaf supplies the nonnegative
literal $a-u$ in its final field.  Theorem~\ref{thm:exact} then compiles
any unrestricted K219 numerical derivation into the same source.

\subsection{Exact packing and cofinal enumeration}

The four command wire formats are
\[
 \begin{aligned}
 f(\MP(z))&=[0,z], &f(\DROP(z))&=[1,z],\\
 f(\NIP(z))&=[2,z],&f(\AX(c,z,y,x))&=[c+3,z,y,x].
 \end{aligned}
\]
Let $E(fs,r)$ be the right fold of $\pr$ over the field list $fs$, ending
in $r$.  Use the tail successor convention
$\enc([])=0$ and $\enc(k::L)=1+E(f(k),\enc(L))$.  A first candidate
is packed without the outer successor.  Every reader starts with zero
scratches and returns a tail no larger than its input.  The initial
decrement on each positive proof list therefore gives strict draining
of every nonhalting candidate, including an arbitrary initial partial
proof list.

If $\DD\vdash1$, let $\mathcal B_1$ be its strong block and put
\[
 R=\enc\bigl(\mathcal B_1[\DROP(0),\NIP(0),\NIP(0)]\bigr).
\]
For each padding value $Y\in\mathbb N$, prepend
$\AX(0,1,1,Y)$.  This AX resets counter and cache to zero from every
incoming retained tuple.  The block then halts or returns top $1$,
and the final three commands halt in the latter case.  The corresponding
candidate number is exactly
\[
 N_Y=\pr(3,\pr(1,\pr(1,\pr(Y,R))))\geq Y.
\]
Thus the family of refuting candidates is cofinal, uniformly in all
incoming top, stack, counter, and cache values.  Choosing one beyond the
initial next-candidate cutoff, and using finite list draining, proves
completeness from every clean source root.  Soundness requires only the
usual theoremhood of top, cache, and stack heads.

\subsection{Arbitrary initial selector and scratch values}

The six-field seed extension can be proved directly for this source.
Write its nine registers as
\[
 s=(N,T,W,A,D,H,X,Y,L),
\]
where $A$ is now the selector register, $D$ the additive counter, $X,Y$
the two scratches, and $L$ the proof list.  Define
\[
 (N',X',n)=
 \begin{cases}
 (N+X+1,0,N),&L=0,\\
 (N,X,L-1),&L>0,
 \end{cases}
 \qquad (u,v)=\unp(\pr(X',Y)+n).
\]
The virtual clean root state is
\[
 s^*=(N',T,W,0,D,H,0,0,\pr(A+u,v)+1).
\]
The unchanged fetch and selector-read prefix takes \emph{both} $s$ and
$s^*$ to exactly
\[
 (N',T,W,A+u,D,H,0,0,v).
\]
Indeed, if $L=0$, the fetch macro copies $N$ into the proof list, adds
it to the existing scratch $X$, transfers that sum back to $N$, and
increments $N$.  Its result is $N+X+1$, scratch zero, and proof list
$N$.  If $L>0$, fetch simply decrements $L$.  The selector reader starts
its unpair walk at $(X',Y)$ and advances by $n$, ending at $(u,v)$.
Its transfer loops add $u$ to the old selector and move $v$ into the
proof list, clearing both scratches.  From $s^*$, the positive proof
list decrements to $\pr(A+u,v)$ and a zero-scratch read returns
$(A+u,v)$, giving the same state.

This agreement occurs \emph{before} the retained common read.  Therefore
both runs execute that read and all later commands identically.
The map preserves the formula data $T,W,H$.  The clean-root calibration
and this identity imply the following exact result.
\begin{corollary}[Common-input source from six arbitrary fields]
For every $N,A,D,X,Y,L\in\mathbb N$, the modified common-input source
started at its genuine root with state
\[
 (N,0,0,A,D,0,X,Y,L)
\]
halts if and only if $\DD\vdash1$, equivalently if and only if
$\CC_{219}\vdash1$.
\end{corollary}
This is a one-time root absorption, not permission to perturb scratches
between commands.  A binary initializer must still reach a complete
represented root frame with the required head position, dispatcher,
register words, header, gap, spare list, and exterior tape.  The source
theorem enlarges the permitted logical seed values; it does not remove
any physical representation obligation.

\subsection{Source provenance and exact proof status}

The source constructor's \texttt{--common} mode reconstructs the
206-distinct-subtree reference source from the accepted K219 source.
It removes the former cache-field read, replaces the former fourth-field
construction by a copy of the old counter, and schedules the new reads
after that copy.  Its default placement is named \texttt{after\_u}.
A separate exact ordinal control-flow comparison checks later
rebracketings of this reference source.  The reference and its calibrated
rebracketing have 4,751 primitive occurrences.

The candidate source recorded for this theorem is
\texttt{source\_0059.json}, with SHA-256
\begin{center}\ttfamily\footnotesize
264fba90229974a53c2c5c943cf5fe8924327b6bbc31bd47e64d7f5c670b7495.
\end{center}
Its reported structural counts are 203 distinct subtrees, maximum depth
45, and maximum zero run 6.  The independent source check establishes
exact ordinal CFG equality with \texttt{a194\_b194\_f1.json}; all 186
original primitive paths outside the AX subtree at address
\texttt{11101001001} preserve their opcode and operand exactly.
The 6,176 guarded decoder cases are checked separately by the source
verification artifacts; their exact receipt is included as
\texttt{evidence/common\_source\_receipt.json}.  These counts and finite checks identify the
source and its control-flow representation; they are not substituted
for the all-natural macro argument or the universal block compiler.

This section is an ordinary mathematical source/calculus proof.
Together with the finite dominated-row certificate, it supplies the
logical and enumeration components of a ZFC-formalizable calibration.
No newly compiled Lean or Rocq endpoint for this modified source is
asserted here.  The exact binary transition table, its all-natural
runtime contracts, and its complete blank-tape startup must be verified
by the separately supplied machine artifacts.
